\documentclass[twocolumn,preprintnumbers,amsmath,amssymb,prl,floatfix]{revtex4-2}
\usepackage{graphicx}
\usepackage{hyperref}
\usepackage{amsmath,amssymb}
\usepackage{algorithm}
\usepackage{algpseudocode}
\usepackage{booktabs}
\usepackage{siunitx}
\usepackage{xcolor}
\hypersetup{
  pdftitle={SCPN Control: Compiling Stochastic Petri Nets into Spiking Neural Network Controllers for Real-Time Tokamak Plasma Control},
  pdfauthor={Miroslav Sotek}
}

\begin{document}

\title{SCPN Control: Compiling Stochastic Petri Nets into Spiking Neural
  Network Controllers for Real-Time Tokamak Plasma Control}

\author{Miroslav Sotek}
\email{protoscience@anulum.li}
\affiliation{ANULUM CH \& LI}


\date{\today}

\begin{abstract}
We present \textsc{scpn-control}, an open-source neuro-symbolic control engine
that compiles Stochastic Petri Nets (SPNs) into spiking neural network (SNN)
controllers with formal contract verification for tokamak plasma control.
The compilation pipeline translates SPN graphs into leaky integrate-and-fire
(LIF) neuron pools with stochastic bitstream encoding, enforcing
pre/post-condition contracts on every control observation and action.
The contribution is this compilation-and-contract architecture and a
runtime-selectable controller stack (PID, nonlinear MPC, $H_\infty$, SNN) sharing
a single contract interface, rather than raw throughput. The control compute is
comfortably within the 4--10~kHz real-time band (native integrated control cycle
\SI{5.62}{\micro\second} P50 on the CI reference runner, reproducible;
Section~\ref{sec:validation}); a fielded loop is bounded by diagnostics,
equilibrium reconstruction, and actuation rather than the controller, so we report
the latency as meeting the budget with margin, not as a competitive speed claim.
The package integrates a multi-layer Kuramoto--Sakaguchi phase dynamics engine
driven by both the Paper~27 $K_{nm}$ coupling matrix~\cite{sotek2026knm} and a
plasma-native 8-layer $K_{nm}$ whose entries encode experimentally grounded
interactions (drift-wave/zonal-flow, NTM/bootstrap-current, ELM/pedestal), with
a Unified Phase Dynamics Equation (UPDE) multi-layer solver monitored by a
Lyapunov sliding-window stability guard.
Five runtime-selectable controllers (PID, MPC, $H_\infty$, SNN,
neuro-cybernetic) operate on a shared Grad--Shafranov equilibrium solver
with L-mode and H-mode profile support.
Disruption prediction with shattered pellet injection (SPI) mitigation,
a Gymnasium-compatible environment, and a digital twin complete the
control stack.
The repository includes provenance-bound equilibrium fixtures and validation
gates, but this manuscript does not promote them to independent experimental or
same-case external-solver validation. The software is distributed under
AGPL-3.0-or-later, with a separate commercial licence available, at
\url{https://github.com/anulum/scpn-control}.
\end{abstract}

\maketitle

%% ════════════════════════════════════════════════════════════════
\section{Introduction}
\label{sec:intro}
%% ════════════════════════════════════════════════════════════════

Real-time plasma control in tokamaks demands sub-millisecond latency,
formal safety guarantees, and runtime adaptability to regime transitions
(L-mode to H-mode, ELM pacing, disruption avoidance).
Existing plasma control systems (PCS) at DIII-D~\cite{ferron1998},
ITER~\cite{iter_codac,humphreys2015}, and EAST rely on classical PID or
model-predictive control with hand-tuned gain schedules.
These approaches lack formal verification of control actions and cannot
adapt their control topology at runtime without operator intervention.

Recent work by Degrave et al.~\cite{degrave2022} demonstrated deep
reinforcement learning for magnetic control on TCV, achieving multi-objective
shaping with a single neural network policy.  While impressive, the approach
treats the controller as a black box---no formal contract checking is
applied to outputs before they reach the plant, and the policy cannot be
inspected for reachability or liveness properties.

Spiking neural networks (SNNs) offer event-driven, low-latency
computation that maps naturally to real-time control
tasks~\cite{tavanaei2019,maass1997}.  Stochastic Petri Nets (SPNs)
provide a formal framework for modeling concurrent, stochastic processes
with well-defined reachability, boundedness, and liveness
properties~\cite{marsan1995,murata1989}.

We bridge these paradigms.  \textsc{scpn-control} compiles an SPN
control specification into an SNN controller, preserving the formal
contract structure of the Petri net as pre/post-condition checks on
every control action.  The SPN graph is declarative and inspectable;
the compiled SNN is fast and neuromorphic-hardware-compatible.

\paragraph{Contributions.}
\begin{enumerate}
  \item A compilation pipeline from SPN graphs to LIF neuron pools
        with stochastic bitstream encoding, antithetic variance
        reduction, and formal contract verification
        (Section~\ref{sec:compilation}).
  \item A multi-layer Kuramoto--Sakaguchi phase dynamics engine with
        the Paper~27 $K_{nm}$ coupling matrix, a plasma-native
        8-layer $K_{nm}$ grounded in experimentally established
        interactions~\cite{diamond2005,terry2000,lahaye2006,porcelli1996},
        UPDE multi-layer solver, PAC gating, and Lyapunov stability
        monitoring (Section~\ref{sec:phase}).
  \item Five runtime-selectable controllers (PID, MPC, $H_\infty$,
        SNN, neuro-cybernetic) sharing a common contract interface,
        with a Gymnasium-compatible environment for RL benchmarking
        (Section~\ref{sec:control}).
  \item A CPU-only Rust-accelerated runtime with a measured
        \SI{5.62}{\micro\second} P50 native integrated control cycle on the
        named CI reference runner; this is controller-side timing, not an
        end-to-end facility-loop claim (Section~\ref{sec:validation}).
  \item A Python package and five-crate Rust workspace with configured
        coverage, documentation, security, and polyglot verification gates
        (Section~\ref{sec:software}).
\end{enumerate}


%% ════════════════════════════════════════════════════════════════
\section{Related Work}
\label{sec:related}
%% ════════════════════════════════════════════════════════════════

\textbf{Classical PCS.}\quad
The DIII-D PCS~\cite{ferron1998} runs real-time equilibrium
reconstruction (EFIT) at 4--10~kHz on dedicated hardware, coupled with
PID shape controllers.  ITER's CODAC~\cite{iter_codac} specifies
sub-millisecond plant I/O but relies on classical control algorithms.
Walker and Humphreys~\cite{walker2020} formalized vertical stability
feedback but without runtime contract verification.

\textbf{ML-based control.}\quad
Degrave et al.~\cite{degrave2022} trained a deep RL policy for
multi-objective TCV shape control.  The policy achieves real-time
execution but provides no formal guarantees on output bounds.
TORAX~\cite{fanni2024torax} provides a differentiable JAX-based
transport simulator suitable for gradient-based optimization but does
not include a control stack.

\textbf{Spiking neural networks.}\quad
Third-generation neural networks~\cite{maass1997} have been proposed
for neuromorphic control but not previously applied to tokamak plasma
control.  Stochastic computing with bitstream
encoding~\cite{alaghi2013} enables low-precision, fault-tolerant
arithmetic suited to noisy plasma environments.

\textbf{Petri nets in control.}\quad
Petri nets have been used for supervisory control and discrete-event
systems~\cite{murata1989}.  We study an SPN-to-SNN compilation path for
continuous plasma-control research while making no priority claim over the
broader neural, supervisory-control, or neuromorphic-control literature.


%% ════════════════════════════════════════════════════════════════
\section{SPN to SNN Compilation}
\label{sec:compilation}
%% ════════════════════════════════════════════════════════════════

\subsection{Stochastic Petri Net Specification}

A Stochastic Petri Net is a tuple
$\mathcal{N} = (P, T, A, W, M_0, \Lambda)$
where $P$ is a finite set of places, $T$ a finite set of transitions,
$A \subseteq (P \times T) \cup (T \times P)$ the arc set,
$W: A \to \mathbb{N}$ the arc weights,
$M_0: P \to \mathbb{N}_0$ the initial marking, and
$\Lambda: T \to \mathbb{R}_{>0}$ the stochastic firing rates.

Places represent physical observables ($I_p$, $W_\mathrm{MHD}$,
$\bar{n}_e$, vertical position $Z$) and controller states.
Transitions encode control decisions (heating power, fueling, vertical
stability corrections).  The SPN is specified declaratively; the user
does not write neural network code.

\subsection{Compilation to LIF Neurons}
\label{sec:lif}

The \texttt{FusionCompiler} maps each transition $t_j \in T$ to a
stochastic leaky integrate-and-fire (LIF) neuron.  The marking
$M(p_i)$ of each place is encoded as a Bernoulli bitstream of length
$L$:
%
\begin{equation}
  b_i^{(k)} \sim \mathrm{Bernoulli}\!\left(\frac{M(p_i)}{M_{\max}}\right),
  \quad k = 1, \ldots, L
  \label{eq:bitstream}
\end{equation}
%
where $L \geq 64$ and $M_{\max}$ normalizes markings to $[0,1]$.
The bitstreams are packed into \texttt{uint64} words
($n_\mathrm{words} = \lceil L/64 \rceil$) with per-element
splitmix64 seed derivation for statistical independence.

Arc weights $W(p_i, t_j)$ define two incidence matrices:
$\mathbf{W}_\mathrm{in} \in \mathbb{R}^{|T| \times |P|}$ (consumption)
and $\mathbf{W}_\mathrm{out} \in \mathbb{R}^{|P| \times |T|}$
(production).  Each LIF neuron integrates:
%
\begin{equation}
  V_j(t+1) = V_j(t) + \frac{\Delta t}{\tau_m}
  \bigl(V_\mathrm{rest} - V_j(t)\bigr)
  + I_j(t) + \sigma\,\eta_j(t)
  \label{eq:lif}
\end{equation}
%
where $I_j = [\mathbf{W}_\mathrm{in}\,\mathbf{m}]_j$ is the
synaptic input from the current marking $\mathbf{m}$,
$\tau_m = 10$~ms, and $\sigma = 0.1$ is the noise amplitude.
When $V_j > V_\mathrm{th}$, the neuron fires ($f_j = 1$) and resets
to $V_\mathrm{reset} = 0$.

The marking update rule preserves the Petri net semantics:
%
\begin{equation}
  \mathbf{m}_{k+1} = \mathrm{clip}\!\left(
    \mathbf{m}_k - \mathbf{W}_\mathrm{in}^\top \mathbf{f}_k
    + \mathbf{W}_\mathrm{out}\,\mathbf{f}_k,\; 0,\; 1\right)
  \label{eq:marking}
\end{equation}

Algorithm~\ref{alg:compile} summarizes the compilation and
Algorithm~\ref{alg:tick} the runtime tick.

\begin{figure}[t]
\hrule\smallskip
\textbf{Algorithm 1:} SPN $\to$ SNN Compilation
\label{alg:compile}
\smallskip\hrule\smallskip
\begin{algorithmic}[1]
\Require SPN $\mathcal{N} = (P, T, A, W, M_0, \Lambda)$,
         bitstream length $L$
\State Build $\mathbf{W}_\mathrm{in}$, $\mathbf{W}_\mathrm{out}$
       from arc set $A$ and weights $W$
\For{each transition $t_j \in T$}
  \State Create LIF neuron with
         $V_\mathrm{th} = \Lambda(t_j)$
\EndFor
\State Pack $\mathbf{W}_\mathrm{in}$ into
       \texttt{uint64} bitstreams (Eq.~\ref{eq:bitstream})
\State Attach contracts from \texttt{ControlObservation},
       \texttt{ControlAction}
\State \Return \texttt{CompiledNet}
       $(\mathbf{W}_\mathrm{in}, \mathbf{W}_\mathrm{out},
         \mathrm{neurons}, \mathrm{contracts})$
\end{algorithmic}
\smallskip\hrule
\end{figure}

\begin{figure}[t]
\hrule\smallskip
\textbf{Algorithm 2:} Control Tick (Runtime)
\label{alg:tick}
\smallskip\hrule\smallskip
\begin{algorithmic}[1]
\Require \texttt{CompiledNet}, observation $\mathbf{o}_t$,
         targets $\mathbf{r}$, marking $\mathbf{m}_t$
\State \textbf{Pre-check:} verify $\mathbf{o}_t$ within
       physical bounds
\State $\mathbf{e}_t \gets (\mathbf{r} - \mathbf{o}_t) /
       \mathbf{s}$ \Comment{normalized error}
\State Inject $\mathbf{e}_t$ into designated places
\State $\mathbf{I} \gets \mathbf{W}_\mathrm{in}\,\mathbf{m}_t$
       \Comment{synaptic input}
\State $\mathbf{f}_t \gets \mathrm{LIF\_step}(\mathbf{I})$
       \Comment{fire vector, Eq.~\ref{eq:lif}}
\State $\mathbf{m}_{t+1} \gets$ Eq.~\ref{eq:marking}
\State $\mathbf{a}_t \gets \mathrm{decode}(\mathbf{m}_{t+1})$
       \Comment{push-pull + slew + clamp}
\State \textbf{Post-check:} verify $\mathbf{a}_t$ within
       actuator limits
\State \Return $\mathbf{a}_t$, $\mathbf{m}_{t+1}$
\end{algorithmic}
\smallskip\hrule
\end{figure}

\subsection{Stochastic Computing Acceleration}

The dense forward pass $\mathbf{y} = \mathbf{W}\,\mathbf{x}$
can be computed via stochastic bitstream arithmetic: input
bitstreams are AND-ed with weight bitstreams, and the output is
the popcount:
%
\begin{equation}
  y_i = \frac{1}{L} \sum_{k=1}^{L}
  \mathrm{popcount}\!\left(
    \mathbf{w}_i^{(k)} \mathbin{\&} \mathbf{x}^{(k)}
  \right)
\end{equation}
%
With \texttt{sc-neurocore} $\geq 3.8.0$, this can use
\texttt{VectorizedSCLayer} with a Rust SIMD backend and packed
\texttt{uint64} operations. Performance is reported only through committed,
host-labelled benchmark artefacts.

\textbf{Antithetic variance reduction.}\quad
For $n$ stochastic passes, we use antithetic sampling:
for each base sample $u \sim \mathrm{Uniform}(0,1)$,
we evaluate both $u < p$ and $u > 1-p$, halving the
variance of the firing probability estimator at no
additional memory cost.

\subsection{Formal Contract Verification}
\label{sec:contracts}

Every control tick verifies two contract layers:

\textbf{Observation contracts.}\quad
The observation $\mathbf{o}_t = (R_\mathrm{axis}, Z_\mathrm{axis})$
is checked against physical bounds.  Errors are normalized and clipped
to $[-1, 1]$ before injection into the SPN marking.

\textbf{Action contracts.}\quad
The decoded action $\mathbf{a}_t = (\Delta I_\mathrm{PF3},
\Delta I_\mathrm{PF,tb})$ is subject to slew-rate limiting
($|\dot{u}| \leq \dot{u}_\mathrm{max}$) and absolute
saturation ($|u| \leq u_\mathrm{abs}$).
Violations trigger a safe fallback (zero-action or last-known-good).

\textbf{Physics invariants.}\quad
Five runtime invariants are monitored:
$q_\mathrm{min} > 1$ (Kruskal--Shafranov~\cite{freidberg2014}),
$\beta_N < 2.8$ (Troyon limit~\cite{troyon1984}),
Greenwald fraction $< 1.2$~\cite{greenwald2002},
$T_i < 25$~keV (first-wall constraint), and
$|\Delta W/W| < 0.01$ (energy conservation).
Violations exceeding 20\% of the threshold are classified as
critical and trigger immediate protective action.


%% ════════════════════════════════════════════════════════════════
\section{Phase Dynamics Engine}
\label{sec:phase}
%% ════════════════════════════════════════════════════════════════

\subsection{Kuramoto--Sakaguchi with Global Field Driver}

Following the Kuramoto model~\cite{kuramoto1975} with the
Sakaguchi phase frustration extension~\cite{sakaguchi1986}
and SCPN Paper~27~\cite{sotek2026knm}, we model $N$ coupled
oscillators across $L = 16$ layers:
%
\begin{equation}
  \dot{\theta}_n^{(\ell)} = \omega_n
  + \frac{1}{N_\ell} \sum_{m=1}^{N_\ell} K_{\ell\ell}
    \sin\!\left(\theta_m^{(\ell)} - \theta_n^{(\ell)} - \alpha\right)
  + \zeta_\ell \sin\!\left(\Psi - \theta_n^{(\ell)}\right)
  \label{eq:kuramoto}
\end{equation}
%
where $\omega_n$ are natural frequencies drawn from the Paper~27
canonical set, $K_{\ell\ell'}$ is the inter-layer coupling matrix,
$\alpha$ is the Sakaguchi phase-lag parameter,
$\zeta_\ell$ is the per-layer global field coupling strength, and
$\Psi$ is the mean-field phase~\cite{acebron2005}.

The order parameter measures synchronization:
%
\begin{equation}
  R_\ell\, e^{i\Psi_\ell}
  = \frac{1}{N_\ell} \sum_{n=1}^{N_\ell} e^{i\theta_n^{(\ell)}}
  \label{eq:order}
\end{equation}
%
$R_\ell = 1$ indicates full synchronization; $R_\ell \to 0$
indicates incoherence.

\subsection{Paper 27 $K_{nm}$ Coupling Matrix}

The $16 \times 16$ coupling matrix encodes inter-layer interactions
with exponential distance decay and empirical calibration anchors:
%
\begin{equation}
  K_{nm} = K_\mathrm{base}\, e^{-\alpha_K |n - m|}
  + \sum_{(i,j) \in \mathcal{A}} c_{ij}\,\delta_{ni}\delta_{mj}
  \label{eq:knm}
\end{equation}
%
with $K_\mathrm{base} = 0.45$, decay $\alpha_K = 0.3$, and
calibration anchors $\mathcal{A}$:
%
\begin{center}
\begin{tabular}{cccc}
\toprule
$(n,m)$ & $(1,2)$ & $(2,3)$ & $(3,4)$ \\
$c_{nm}$ & $0.302$ & $0.201$ & $0.252$ \\
\midrule
$(n,m)$ & $(4,5)$ & $(1,16)$ & $(5,7)$ \\
$c_{nm}$ & $0.154$ & $\geq 0.05$ & $\geq 0.15$ \\
\bottomrule
\end{tabular}
\end{center}
%
The last two entries are cross-hierarchy boosts (L1$\leftrightarrow$L16
and L5$\leftrightarrow$L7) from Paper~27 \S4.3.

\subsection{Plasma-Native $K_{nm}$ Coupling Matrix}
\label{sec:plasma_knm}

The Paper~27 matrix was derived from a 16-layer model of consciousness
dynamics.  Tokamak plasmas have their own multi-scale hierarchy.  We
construct a dedicated $8 \times 8$ coupling matrix whose entries
map directly onto experimentally established plasma physics
interactions.

\textbf{Layer hierarchy.}\quad
Eight layers span the full range of tokamak timescales from
micro-turbulence (${\sim}\mu$s) to plasma--wall interaction
(${\sim}$s--min):
%
\begin{center}
\begin{tabular}{clc}
\toprule
Layer & Process & Timescale \\
\midrule
P0 & Micro-turbulence (ITG/TEM) & $\mu$s \\
P1 & Zonal flows ($E{\times}B$ shear) & ms \\
P2 & MHD / tearing (NTM, 2/1 island) & ms--s \\
P3 & Sawtooth / ELM cycles & 10--100~ms \\
P4 & Transport barrier (H-mode pedestal) & $\sim$100~ms \\
P5 & Current profile ($q(\rho), l_i$) & s \\
P6 & Global equilibrium (GS $\psi(R,Z)$) & s \\
P7 & Plasma--wall interaction & s--min \\
\bottomrule
\end{tabular}
\end{center}

\textbf{Coupling construction.}\quad
The backbone follows exponential distance decay as in
Eq.~\eqref{eq:knm}, but with reduced base strength
$K_\mathrm{base} = 0.30$ and steeper decay $\alpha_K = 0.5$
to reflect the wider timescale separation in plasma:
%
\begin{equation}
  K_{nm}^\mathrm{plasma} = K_\mathrm{base}\, e^{-\alpha_K |n-m|}
  + \sum_{(i,j)\in\mathcal{P}} p_{ij}\,\delta_{ni}\delta_{mj}
  \label{eq:plasma_knm}
\end{equation}
%
The physics overlays $\mathcal{P}$ encode seven experimentally
grounded interactions:
%
\begin{center}
\begin{tabular}{cll}
\toprule
$(n,m)$ & $p_{nm}$ & Reference \\
\midrule
(0,1) & 0.42 & Drift-wave/zonal predator--prey~\cite{diamond2005} \\
(1,4) & 0.28 & $E{\times}B$ shear suppression~\cite{terry2000} \\
(2,5) & 0.35 & NTM $\leftrightarrow$ bootstrap current~\cite{lahaye2006} \\
(3,5) & 0.30 & Sawtooth $\leftrightarrow$ current redistribution~\cite{porcelli1996} \\
(3,4) & 0.32 & ELM crash depletes pedestal \\
(4,6) & 0.25 & Transport $\leftrightarrow$ GS equilibrium~\cite{lutjens1996} \\
(7,4) & 0.20 & PWI recycling modifies edge transport~\cite{stangeby2000} \\
\bottomrule
\end{tabular}
\end{center}
%
All couplings are symmetrised ($K_{nm} = K_{mn}$) and clamped
non-negative after overlay application.

\textbf{Instability mode bias.}\quad
A mode parameter $m \in \{\texttt{baseline}, \texttt{elm},
\texttt{ntm}, \texttt{sawtooth}, \texttt{hybrid}\}$ multiplicatively
amplifies couplings associated with the dominant instability.
For example, \texttt{elm} mode scales $K_{34}$ and $K_{43}$ by
$1.8\times$ (stronger pedestal--ELM interaction), while
\texttt{hybrid} mode elevates all entries by $1.15\times$ to
represent the tighter multi-scale coupling of advanced scenarios.
This allows the phase engine to reflect the qualitative
synchronization topology of a given operating regime without
retuning individual entries.

\textbf{Machine-parameter scaling.}\quad
A configuration-driven builder adjusts $K_\mathrm{base}$ via a
dimensionless beta proxy:
%
\begin{equation}
  K_\mathrm{base}(R_0, a, B_0, n_e)
  = 0.30\,\bigl(1 + 0.5\,\tilde{\beta}\bigr),\quad
  \tilde{\beta} = \frac{n_e\, a}{B_0^2}
  \label{eq:kbase_beta}
\end{equation}
%
where $n_e$ is in units of $10^{19}$~m$^{-3}$, $a$ in metres, and
$B_0$ in Tesla.  Higher beta implies more free energy driving
instabilities, hence stronger multi-scale coupling.  The
cylindrical safety factor $q_\mathrm{cyl} = 5 a^2 B_0 / (R_0 I_p)$
is computed simultaneously; when $q_\mathrm{cyl} < 1$ the builder
automatically selects sawtooth mode, reflecting the physical
inevitability of internal kink activity at low $q$.

\textbf{Design for online adaptation.}\quad
The static mode-switchable architecture is designed as the
offline-verified foundation for a future real-time adaptive $K_{nm}$
where diagnostic signals update coupling strengths online.
Mirnov coil amplitudes would modulate $K_{27}$ (MHD--PWI),
Thomson scattering pedestal gradients would drive $K_{34}$
(ELM--transport), and soft X-ray sawtooth inversion radius would
scale $K_{35}$ (sawtooth--current).
The mode-switching interface already supports per-tick calls with
${\sim}15\,\mu$s overhead; the remaining step is closing the loop
from diagnostic pre-processing to coupling update, which we address
in Section~\ref{sec:conclusion}.

\subsection{UPDE Multi-Layer Solver}

The Unified Phase Dynamics Equation (UPDE) couples all 16 layers.
For oscillator $n$ in layer $\ell$:
%
\begin{multline}
  \dot{\theta}_n^{(\ell)} = \omega_n^{(\ell)}
  + \frac{1}{N_\ell} \sum_m K_{\ell\ell}\,
    \sin\!\left(\theta_m^{(\ell)} - \theta_n^{(\ell)}\right)
  \\
  + \sum_{\ell' \neq \ell} K_{\ell\ell'}\, R_{\ell'}\,
    \sin\!\left(\Psi_{\ell'} - \theta_n^{(\ell)}\right)
  + \zeta_\ell \sin\!\left(\Psi - \theta_n^{(\ell)}\right)
  \label{eq:upde}
\end{multline}

\textbf{PAC gating.}\quad
Optional phase-amplitude coupling modulates cross-layer
drive based on source coherence:
%
\begin{equation}
  \tilde{K}_{\ell\ell'} = K_{\ell\ell'}
  \left[1 + \gamma_\mathrm{PAC}\,(1 - R_{\ell'})\right]
\end{equation}
%
When source layer $\ell'$ desynchronizes ($R_{\ell'} \downarrow$),
inter-layer coupling increases, implementing a stabilizing
negative feedback. The repository benchmark harness measures the added compute
path, but no portable PAC-overhead value is admitted in this manuscript.

\subsection{Lyapunov Stability Guard}
\label{sec:lyapunov}

A sliding-window monitor tracks the Lyapunov candidate
function~\cite{acebron2005}:
%
\begin{equation}
  V(t) = \frac{1}{N}\sum_{n=1}^{N}
  \left(1 - \cos(\theta_n - \Psi)\right)
\end{equation}
%
and computes the exponential rate over a configurable window
of $W = 50$ samples:
%
\begin{equation}
  \lambda = \frac{1}{\Delta t}\,
  \ln\!\left(\frac{V(t)}{V(t - W\Delta t)}\right)
  \label{eq:lambda}
\end{equation}
%
$\lambda < 0$ indicates convergence toward synchronization.
If $\lambda > 0$ for $K_\mathrm{max} = 3$ consecutive windows,
the guard triggers a coherence recovery action (increased $\zeta$
or controller gain boost).

Table~\ref{tab:lyapunov} shows $\lambda$ as a function of global
field strength $\zeta$ for $N = 1000$ oscillators.

\begin{table}[tb]
\centering
\caption{Lyapunov exponent $\lambda$ vs.\ global field strength
$\zeta$ ($N = 1000$, 200 steps, $\Delta t = 1$~ms).}
\label{tab:lyapunov}
\begin{tabular}{rcc}
\toprule
$\zeta$ & $\lambda$ ($K=0$) & $\lambda$ ($K=2$) \\
\midrule
$0.0$ & $+0.01$ & $+0.04$ \\
$0.1$ & $-0.03$ & $-0.02$ \\
$0.5$ & $-0.23$ & $-0.24$ \\
$1.0$ & $-0.49$ & $-0.53$ \\
$3.0$ & $-1.65$ & $-1.83$ \\
$5.0$ & $-3.01$ & $-3.35$ \\
\bottomrule
\end{tabular}
\end{table}


%% ════════════════════════════════════════════════════════════════
\section{Control Architecture}
\label{sec:control}
%% ════════════════════════════════════════════════════════════════

\subsection{Controller Suite}

All five controllers share the contract interface of
Section~\ref{sec:contracts}: observation pre-check, action
post-check, and physics invariant monitoring.

\textbf{PID.}\quad
Dual-channel $I_p$ + vertical position control with anti-windup
and derivative low-pass filtering.

\textbf{MPC.}\quad
Gradient-descent optimization over a 10-step prediction horizon
with $L_2$ regularization ($\lambda = 0.1$):
%
\begin{equation}
  J = \sum_{t=0}^{N_p-1}
  \left\|\mathbf{A}\mathbf{x}_t + \mathbf{B}\mathbf{u}_t
  - \mathbf{x}_\mathrm{ref}\right\|^2
  + \lambda\,\|\mathbf{u}_t\|^2
\end{equation}
%
The plant model is a linearized 4-state system
$(\mathbf{x} = [R_\mathrm{axis}, Z_\mathrm{axis},
R_x, Z_x]^\top)$ with coil-current inputs.
Action limits are $\pm 2.0$~A per coil per step.

\textbf{$H_\infty$.}\quad
$\gamma$-iteration robust synthesis~\cite{zhou1996,glover1988}
on a linearized vertical stability plant:
%
\begin{align}
  \dot{\mathbf{x}} &= \mathbf{A}\mathbf{x}
  + \mathbf{B}_1\mathbf{w} + \mathbf{B}_2\mathbf{u}
  \label{eq:hinf_plant} \\
  \mathbf{z} &= \mathbf{C}_1\mathbf{x}
  + \mathbf{D}_{12}\mathbf{u}
  \label{eq:hinf_perf}
\end{align}
%
The controller minimizes $\|\mathbf{T}_{zw}\|_\infty < \gamma$
via bisection on $\gamma$ (range $[1.01, 10^6]$,
$r_\mathrm{tol} = 10^{-3}$, max 100 iterations).
State-feedback and observer gains are obtained from the
continuous algebraic Riccati equations; discrete-time
adaptation uses zero-order-hold discretization.
Robustness margin: $\pm 20\%$ multiplicative plant uncertainty.

\textbf{SNN.}\quad
Compiled from SPN specification
(Section~\ref{sec:compilation}).

\textbf{Neuro-cybernetic.}\quad
Dual radial ($R$) + vertical ($Z$) spiking neuron pools
($n = 20$ neurons each) with push-pull decoding:
%
\begin{equation}
  u = g \cdot \left(\frac{N_+}{N_w} - \frac{N_-}{N_w}\right)
\end{equation}
%
where $N_\pm$ are spike counts in the positive/negative
populations over a window of $N_w$ ticks and $g$ is the
output gain.

\subsection{Grad--Shafranov Equilibrium Solver}

The GS equation in toroidal coordinates:
%
\begin{equation}
  \Delta^* \psi \equiv R\frac{\partial}{\partial R}
  \left(\frac{1}{R}\frac{\partial\psi}{\partial R}\right)
  + \frac{\partial^2\psi}{\partial Z^2}
  = -\mu_0 R^2 \frac{dp}{d\psi}
  - \frac{1}{2}\frac{dF^2}{d\psi}
  \label{eq:gs}
\end{equation}
%
is solved by Picard iteration on an $N_R \times N_Z$ grid
(default $65 \times 65$).

\textbf{Vacuum field.}\quad
For each coil at $(R_c, Z_c)$ with current $I_c$:
%
\begin{equation}
  \psi_\mathrm{vac} = \frac{\mu_0 I_c}{2\pi}
  \sqrt{(R+R_c)^2 + \Delta Z^2}\;
  \frac{(2-k^2)K(k) - 2E(k)}{k^2}
\end{equation}
%
where $k^2 = 4RR_c / [(R+R_c)^2 + \Delta Z^2]$ and
$K$, $E$ are complete elliptic integrals.

\textbf{Profile modes.}\quad
L-mode: $p(\rho) = p_0(1 - \rho^{\alpha_c})$.
H-mode: core profile with mtanh pedestal transition at
$\rho_\mathrm{ped}$ with configurable height and width.

Convergence to $\|\delta\psi\|_\infty < 10^{-6}$ typically
in 12--18 iterations.  Solver backends: Python (SOR, Picard),
Rust (multigrid + SOR), or optional C++ HPC bridge.

\subsection{Disruption Prediction and SPI Mitigation}
\label{sec:disruption}

An 11-dimensional feature vector is extracted from the plasma
state at each tick:
mean, std, max, and slope of the tearing mode signal,
energy ($\langle s^2 \rangle$), last value,
toroidal mode amplitudes ($n = 1, 2, 3$),
toroidal asymmetry index, and radial spread.

The tearing mode physics follows the modified Rutherford
equation~\cite{devries2011}:
%
\begin{equation}
  \frac{dw}{dt} = \Delta'(1 - w/w_\mathrm{sat})
\end{equation}
%
Disruption is triggered when the island width exceeds
$w > 8.0$ (mode lock).

An LSTM classifier (PyTorch, optional) predicts $P(\mathrm{disrupt})$.
When $P > 0.8$, shattered pellet injection
(SPI)~\cite{lehnen2015} is triggered with halo current and
runaway electron post-disruption physics.

The predictor is trained on synthetic data derived from 11
reference DIII-D shot profiles~\cite{rea2019}
(6 anomalous: locked mode, density limit, VDE, tearing, beta
limit; 5 safe: hybrid, H-mode, snowflake, negative delta,
high beta).

\subsection{Digital Twin and Gymnasium Environment}

\textbf{TokamakDigitalTwin.}\quad
A 2D diffusive-reaction model on a $40 \times 40$ poloidal grid
with core Gaussian heating ($+5.0$~K/step), configurable
turbulent diffusivity ($D_\mathrm{turb} = 0.5$ in island
regions), and radiative losses ($\propto T^2$).
Rational surface detection at $q = 1.5, 2.0, 2.5, 3.0$ marks
regions of enhanced transport.

\textbf{TokamakFlightSim.}\quad
Isoflux flight simulator with first-order actuator dynamics
($\tau = 60$~ms), rate limits
($\leq 10^6$~A/s, ITER PF coil spec), sensor noise, and
configurable measurement delay.

\textbf{TokamakEnv.}\quad
Gymnasium-compatible~\cite{gymnasium2023} environment wrapping
the digital twin.  Observation space: 4D continuous
$(R, Z, R_x, Z_x)$.  Action space: continuous coil currents.
Reward: $-\|\mathbf{x} - \mathbf{x}_\mathrm{ref}\|^2$
with disruption penalty.


%% ════════════════════════════════════════════════════════════════
\section{Benchmarks and Validation}
\label{sec:validation}
%% ════════════════════════════════════════════════════════════════

\subsection{Control Loop Timing}

Table~\ref{tab:timing} reports median (P50) and tail (P99) per-step latencies,
measured by the committed reproducible benchmarks (warm-up plus repeated samples,
not single-run averages) on the CI reference runner (AMD EPYC 7763). Every
available backend is reported; a backend whose dependency is absent is listed as
such rather than estimated. The repository publishes a separately measured
Intel i5-11600K workstation column alongside the CI result, without treating
either host as hardware-independent performance.

\begin{table}[tb]
\centering
\caption{Control-loop latency on the recorded AMD EPYC 7763 CI runner. Isolated
controllers use 200 warm-up steps and 2000 measured steps, except MPC with 500
measured steps. The integrated row uses 5000 steps in each of seven repeats.}
\label{tab:timing}
\begin{tabular}{llrrr}
\toprule
Controller & Backend & P50 (\si{\micro\second}) & P99 (\si{\micro\second})
  & Throughput \\
\midrule
PID          & NumPy & 0.42  & 0.52  & 2381~kHz \\
SNN          & NumPy & 23.90 & 43.14 & 42~kHz  \\
SNN          & Rust  & 0.88  & 0.99  & 1136~kHz \\
MPC ($N_p{=}10$) & acados & 10679 & 11393 & 0.09~kHz \\
MPC ($N_p{=}10$) & osqp & 16846 & 18090 & 0.06~kHz \\
MPC ($N_p{=}10$) & internal QP & 18847 & 25172 & 0.05~kHz \\
$H_\infty$   & NumPy (general SS) & 29.43 & 53.42 & 34~kHz \\
$H_\infty$   & Rust (2-state VDE) & 1.39 & 1.49 & 719~kHz \\
\midrule
Integrated control cycle & Rust native & 5.62 & 6.11 & 178~kHz \\
\bottomrule
\end{tabular}
\end{table}

The SNN, $H_\infty$, PID, and native integrated control cycle all sit well within
the 4--10~kHz real-time band on the control-compute side. This is a
necessary-not-sufficient property: a fielded loop is bounded by diagnostics,
equilibrium reconstruction, and actuation, not by these control-compute figures,
so they are reported as meeting the budget rather than as an end-to-end speed
advantage over fielded systems.
The nonlinear MPC is millisecond-scale on every QP backend, with \texttt{acados}
the fastest ($\approx$\SI{10.7}{\milli\second}); the full five-backend comparison
(internal, scipy, osqp, casadi, acados) is regenerated by the repository
benchmark. The microsecond MPC regime is not reachable through the Python
\texttt{step\_rti} path with any QP backend, because per-stage set/get marshalling
across the Python--C boundary dominates the tick; a microsecond MPC tick would
require a non-Python (C/Rust) control loop, not merely a compiled QP solver. The
MPC row is therefore reported as a research-grade tick latency, not a real-time
claim.

\subsection{Capability Context}

The related systems solve different problem slices. DIII-D PCS
demonstrates fielded equilibrium reconstruction and control~\cite{ferron1998};
TORAX provides differentiable transport~\cite{fanni2024torax}; FreeGS provides
an open Grad--Shafranov solver~\cite{freegs}; and the TCV work demonstrates
learned magnetic control~\cite{degrave2022}. The contribution evaluated here is
the repository's explicit SPN-to-SNN compilation and contract boundary. No
overall superiority claim follows from this capability difference.

\subsection{Kuramoto Phase Sync: Python vs.\ Rust}

The repository includes a side-by-side Python/Rust Kuramoto benchmark over
increasing oscillator counts. Those observations are useful for local
regression, but this paper does not reproduce them as a portable speedup claim
because the legacy table is not accompanied by a sealed raw host-context
artefact at the manuscript revision.

\subsection{Equilibrium RMSE}

The repository carries checksum-bound GEQDSK fixtures and regression gates.
These are software-regression evidence, not independent reconstruction accuracy
against a current experimental reference code on identical measurements. The
legacy percentage table is therefore withdrawn from this review draft pending a
same-case, provenance-complete comparison.

\subsection{Phase Dynamics Convergence}

The repository provides deterministic phase-synchronisation scenarios and a
sliding-window Lyapunov monitor. Their outputs establish reproducible simulated
behaviour for the configured model only; they do not establish experimental
plasma stability or an independent fail-closed safety interlock.


%% ════════════════════════════════════════════════════════════════
\section{Software Architecture}
\label{sec:software}
%% ════════════════════════════════════════════════════════════════

\subsection{Python Package}

The Python package is organised by responsibility into \texttt{scpn/},
\texttt{core/}, \texttt{control/}, and \texttt{phase/}, with validation and
public interfaces outside those implementation packages. Exact module and test
inventories are generated from the repository tree rather than frozen in the
manuscript because they change as the software evolves.

\subsection{Rust Workspace}
\label{sec:rust}

Five crates compiled via PyO3~\cite{pyo3}:

\begin{enumerate}
\item \textbf{control-types}: \texttt{PlasmaState},
  \texttt{EquilibriumConfig}, \texttt{ControlAction},
  physical constants.
\item \textbf{control-math} (14 modules): multigrid, SOR,
  GMRES, Chebyshev, FFT, Boris pusher, LIF neuron, Kuramoto
  step, symplectic integrator, tridiagonal solver.
\item \textbf{control-core} (20 modules): GS kernel,
  transport, vacuum field, X-point detection, inverse solver,
  pedestal, stability, RF heating, VMEC/BOUT++ interfaces.
\item \textbf{control-control} (10 modules): PID, MPC,
  $H_\infty$, SNN pool, SPI mitigation, digital twin,
  optimal control, SOC learning.
\item \textbf{control-python}: PyO3 bindings exporting
  \texttt{PyFusionKernel}, \texttt{PySnnPool},
  \texttt{PyMpcController}, \texttt{PyPlasma2D},
  \texttt{PyTransportSolver}, \texttt{PyRealtimeMonitor},
  and SCPN kernels (\texttt{dense\_activations},
  \texttt{marking\_update}, \texttt{sample\_firing}).
\end{enumerate}

Python code auto-detects the Rust backend at import time
via \texttt{importlib.util.find\_spec} and dispatches
hot-path calls accordingly, with a pure-Python fallback
for all functionality.

\subsection{Testing and CI}

The repository runs module-specific Python tests, Rust workspace tests, and
integration checks across its hosted workflows. The configured Python package
coverage gate is 100\%; the manuscript does not substitute a historical test or
workflow count for exact-head CI evidence:
\begin{itemize}
  \item \textbf{Python}: pytest, ruff, mypy, bandit, coverage
        (100\% configured package gate), scientific regression gates
  \item \textbf{Rust}: \texttt{cargo test}, clippy, fmt, audit,
        Criterion benchmarks (Kuramoto, transport, LIF, Boris)
  \item \textbf{Integration}: named public/runtime paths, documentation builds,
        and generated-evidence drift checks
\end{itemize}


%% ════════════════════════════════════════════════════════════════
\section{Limitations}
\label{sec:limitations}
%% ════════════════════════════════════════════════════════════════

We state limitations explicitly so they are not overstated
by inference:

\begin{enumerate}
  \item \textbf{Equilibrium}: fixed- and free-boundary implementations are
    repository software surfaces; independent same-case reconstruction against
    a current external reference and facility diagnostics remains open.
  \item \textbf{Transport}: native and adapter surfaces do not establish
    quantitative equivalence to TGLF, GENE, GS2, CGYRO, or QuaLiKiz without
    admitted external artefacts on identical inputs. Nonlinear CBC saturation
    and unit-normalised cross-code agreement remain open.
  \item \textbf{Validation}: repository fixtures include synthetic and public
    reference data. Offline replay is not live MDSplus acquisition or a facility
    experiment, and no fielded PCS validation is claimed.
  \item \textbf{Disruption prediction}: available measured-shot evidence is
    offline and dataset-bounded; it is not a deployed predictor or actuator
    authority.
  \item \textbf{Deployment}: CODAC, EPICS, IMAS, HDL-export, and hardware-facing
    adapters are integration surfaces, not proof of installation, post-route
    timing, HIL operation, or certified plant control.
  \item \textbf{Plasma $K_{nm}$ calibration}: Coupling strengths
    are order-of-magnitude estimates derived from published
    scalings, not fitted to experimental multi-diagnostic
    correlation data.  Real-time adaptive mode requires
    closing the diagnostic-to-coupling loop.
  \item \textbf{Rust acceleration}: optional acceleration is reported only
    where paired, host-labelled artefacts exist. Performance does not imply
    mathematical parity outside the explicitly tested path.
\end{enumerate}


%% ════════════════════════════════════════════════════════════════
\section{Conclusion and Future Work}
\label{sec:conclusion}
%% ════════════════════════════════════════════════════════════════

We presented a neuro-symbolic approach to tokamak plasma control
that compiles Stochastic Petri Nets into spiking neural network
controllers through a formal pipeline preserving contract
verification at every control tick.

The technical result is an inspectable compilation-and-contract architecture:
control logic is specified as a Petri-net graph, compiled into LIF neuron pools,
and checked at the observation/action boundary. The committed host-labelled
benchmarks show microsecond-scale controller-side native handoff on their named
machines, but do not compare end-to-end facility loops. The phase engine adds a
plasma-native 8-layer $K_{nm}$ grounded in
drift-wave/zonal-flow~\cite{diamond2005},
NTM/bootstrap-current~\cite{lahaye2006}, and
ELM/pedestal interactions, with machine-parameter scaling
via $\tilde{\beta}$ and automatic sawtooth-mode selection
at $q_\mathrm{cyl} < 1$.

Future work targets:
\begin{itemize}
  \item \textbf{Real-time adaptive $K_{nm}$}: closing the loop
        from diagnostic pre-processing (Mirnov coils, Thomson
        scattering, soft X-ray) to per-tick coupling updates,
        making $K_{nm}(t)$ a function of the instantaneous
        plasma state rather than a static operating-point selection
  \item Same-case independent equilibrium and transport-code comparison
  \item Shot-disjoint measured-data validation with uncertainty and failure accounting
  \item Named-target post-route and hardware-in-the-loop evidence
  \item Facility-approved non-actuating shadow operation
\end{itemize}

The software is available at \url{https://github.com/anulum/scpn-control}
under AGPL-3.0-or-later, with a separate commercial licence available.

%% ════════════════════════════════════════════════════════════════
\begin{acknowledgments}
This work builds on the SCPN theoretical framework (Papers 1--27)
and the \texttt{scpn-fusion-core} codebase.
We thank the open-source fusion community for reference equilibria,
the DIII-D team for published GEQDSK data, and the Commonwealth
Fusion Systems team for SPARC equilibrium parameters.
\end{acknowledgments}

\bibliography{references}
\bibliographystyle{apsrev4-2}

\end{document}
